% Define block styles 
\tikzstyle{decision} = [diamond, draw, fill=blue!20, 
text width=4em, text badly centered, node distance=2cm, inner sep=1pt] 
\tikzstyle{block} = [rectangle, draw, fill=blue!20, 
text width=3cm, text centered, rounded corners, minimum height=2em] 
\tikzstyle{line} = [draw, -latex'] 
\tikzstyle{cloud} = [draw, ellipse,fill=red!20, node distance=4cm, 
minimum height=4em]




%\pagestyle{empty} 
\thispagestyle{empty}
\begin{frame}[label=Mapa]
\begin{tikzpicture}[node distance = 1.5cm, auto] 
           %%%% Mapa de nodos%%%%


\node [block] (Titulo) at (6,6) {Convergencia en ENO};

\node [cloud] (Orden) at (3,4) { \hyperlink{Orden}{ $x-$convergencia }}; 
\node [cloud] (Norma) at (9,4) {\hyperlink{Norma}{Convergencia  en norma}}; 
\node [block] (EB) at (6,2) {\hyperlink{EBO}{Espacios de normados Ordenados} } ; 
%\node [block, below of=EB, node distance=2cm] (comparar) {\hyperlink{Comparacion}{Comparación}}; 
% Dibujando las flechas 
\path [line] (Titulo) --  (Orden); 
%\path [line] (Titulo) --  (Cono); 
\path [line] (Titulo) -- (Norma); 
\path [line,dashed] (Orden) -- (EB);

%\path [line,dashed] (Cono) --  (Orden); 
%\path [line,dashed] (Cono) -- (Norma); 
\path [line,dashed] (Norma) -- (EB); 
%	\path [line] (EB) -- (comparar); 
\end{tikzpicture}	
\end{frame}
